\section{Why models pass the ceiling: document support}\label{sec:support}

Section~\ref{sec:results} reports task--split combinations on which models score below the estimated floor ($f>1$). On the 17 tasks of the preprint these looked like isolated anomalies. With 41 tasks they are not isolated: ECFP with gradient boosting exceeds the floor on 16 of 41 tasks under a random split. This section identifies the cause, and the cause turns out to say more about the splits than about the floor.

\paragraph{$f$ scales with the noise it is meant to divide out.} The fraction of ceiling was introduced to put tasks with different noise levels on one scale. It does not. Across the 41 tasks, mean $f$ of ECFP+GBM on random splits rises with the estimated noise $\hat s$ (Spearman $\rho=+0.74$, $p=3\times10^{-8}$). This is not an artefact of the vanishing headroom in the denominator of $f$: the ratio $\mathrm{RMSE}_{\mathrm{floor}}/\mathrm{RMSE}$, which has no such denominator, rises with $\hat s$ just as strongly ($\rho=+0.80$, $p=4\times10^{-10}$). On noisy tasks models really do get closer to, and past, the floor estimated for them.

\paragraph{The cause is how much of a test molecule's document is in training.} For every test molecule measured in a single document --- 609{,}539 single-document test cases over the 41 tasks --- we count how many molecules from that same document are in the training set (\texttt{scripts/}\allowbreak\texttt{14\_document\_support.py}). Test error falls monotonically with that count, and crosses the floor (Table~\ref{tab:support}).

\begin{table}[t]
\centering\small
\caption{Test RMSE (pChEMBL) of ECFP+GBM against the number of training molecules from the test molecule's own document, over 41 tasks, 5 seeds, test molecules measured in one document. The median $\hat s$ over these tasks is 0.54. Counts are molecule--split--seed cases.}
\label{tab:support}
\begin{tabular}{lrrrrrr}
\toprule
& \multicolumn{2}{c}{Random} & \multicolumn{2}{c}{Scaffold} & \multicolumn{2}{c}{Temporal} \\
\cmidrule(lr){2-3}\cmidrule(lr){4-5}\cmidrule(lr){6-7}
Same-document training mols & RMSE & $n$ & RMSE & $n$ & RMSE & $n$ \\
\midrule
0        & 1.073 & 2{,}484  & 1.078 & 8{,}251  & 1.055 & 107{,}050 \\
1--4     & 0.835 & 10{,}862 & 0.894 & 19{,}380 & 1.040 & 42{,}135 \\
5--19    & 0.665 & 58{,}597 & 0.762 & 58{,}551 & 0.821 & 19{,}085 \\
20--99   & 0.593 & 57{,}583 & 0.665 & 49{,}600 & 0.658 & 17{,}145 \\
100+     & 0.506 & 71{,}912 & 0.564 & 65{,}744 & 0.503 & 21{,}160 \\
\bottomrule
\end{tabular}
\end{table}

A model given no molecule from the test molecule's document errs by about 1.05 pChEMBL units, well above the floor; given a hundred or more, it errs by about 0.51, below it. Document support, not chemistry, decides whether a model appears to beat the noise ceiling.

\paragraph{The three splits do not differ in difficulty.} Read the first row of Table~\ref{tab:support} across: with no same-document support the three splits give 1.073, 1.078 and 1.055. They are the same task. What differs is how many test molecules are left in that state: 2{,}484 under a random split against 107{,}050 under a temporal one. The familiar ordering --- random easy, scaffold harder, temporal hardest --- is, on these data, mostly a statement about document support rather than about generalising to new chemotypes or to the future. This reframes the gap reported in Section~\ref{sec:results}. Whether it reframes the similar gaps reported elsewhere in the potency-prediction literature we cannot say from one data set, but the mechanism is not peculiar to our tasks, and the check is cheap for anyone to run.

\paragraph{It is support, not document size.} Support and document size are nearly collinear under a random split: a large document contributes many training molecules by construction. Scaffold splits break the collinearity, because molecules of one document are separated by scaffold. Within each document-size stratum the drop from a barely-seen document to a well-seen one survives, although it is not monotone across the upper cells (Table~\ref{tab:strata}).

\begin{table}[ht]
\centering\small
\caption{Test RMSE of ECFP+GBM by document size and by the share of that document seen in training, scaffold split. Support and document size are nearly collinear under a random split; the scaffold split separates them.}
\label{tab:strata}
\begin{tabular}{lrrrr}
\toprule
Document size & $<$20\% & 20--50\% & 50--80\% & $>$80\% \\
\midrule
1--49 molecules   & 1.004 & 0.812 & 0.754 & 0.715 \\
50--199 molecules & 0.743 & 0.650 & 0.629 & 0.612 \\
200+ molecules    & 0.999 & 0.505 & 0.567 & 0.543 \\
\bottomrule
\end{tabular}
\end{table}

\paragraph{Paired test.} Finally, the same molecules, compared with and without support. We take the 2{,}096 test molecules that belong to documents of at least 100 molecules, that have at least 100 same-document training molecules under the random split, and that have none under the temporal split. Their RMSE is 0.459 with support and 0.913 without (Wilcoxon signed-rank on squared errors, $p=8\times10^{-179}$). Intrinsic difficulty is held fixed by construction; only support changes. The comparison does also move the training data earlier in time, so part of the difference may be genuine temporal shift; the within-split gradients above carry no such confound.

\paragraph{Mechanisms we tested and rejected.} Three plausible explanations do not survive contact with the data, and we record them so that others need not repeat them. (i) \emph{Concentration of a task in few documents}: rank correlation of mean $f$ with the number of documents per task is $-0.04$, and with the share of measurements from the largest document $+0.06$. (ii) \emph{Between-laboratory offsets}: using the author metadata of Section~\ref{sec:provenance}, the per-task gap between $\hat\sigma$ among document pairs with no shared author and pairs that share one does not predict $f$ ($\rho=+0.02$, $n=23$ tasks with at least 20 shared-author pairs). (iii) \emph{Potency-dependent heteroscedasticity}: within replicated molecules, the absolute between-document deviation does not grow with potency (median per-task $\rho=-0.01$; 20 of 41 tasks positive, sign test $p=1.0$). The effect is specific to how much of a test molecule's own document the model has already seen.

\paragraph{What this means for the floor.} The floor itself is not refuted. It estimates the disagreement between independent documents, and Section~\ref{sec:provenance} shows that estimate is if anything conservative, because replicates sharing authors agree better than independent ones. What the comparison shows is that a model evaluated on a random or scaffold split is not solving the task the floor describes: it is partly recovering the offset of a document it has already seen, which an independent measurement of the same molecule would not reproduce. The noise ceiling is therefore most useful not as a normalising constant but as a detector: scores that approach or pass it are evidence of document support, not of models nearing the limit of the data. Two papers reporting a ``scaffold split'' on the same target can differ by a factor of two in error for this reason alone, so the distribution of document support over the test set says more about what a number means than the name of the split does; it is computed from the document identifier ChEMBL already supplies. Section~\ref{sec:docsplit} tests the remedy.
